\documentclass[10pt,a4paper]{amsart}
\usepackage[margin=19mm]{geometry}
\usepackage{amsmath,amssymb,mathtools}
\usepackage[scr=rsfs,cal=euler]{mathalfa}
\usepackage{xfrac}
\usepackage[unicode,hidelinks,bookmarks=false]{hyperref}
\newcommand{\ie}{i.e.\ }
\newcommand{\cf}{cf.\ }
\newcommand{\resp}{respectively\ }
\newcommand{\Subsection}{\subsection}
\providecommand{\nobreakdash}{\nobreak\hbox{-}}
\setlength{\emergencystretch}{2em}
\multlinegap0pt
\usepackage[shortlabels]{enumitem}
\usepackage{amssymb}
\usepackage{mathtools}
\newtheorem{theorem}{Theorem}
\newtheorem{remark}{Remark}
\usepackage{cleveref}
\crefname{theorem}{Theorem}{Theorems}
\crefname{remark}{Remark}{Remarks}
\newcommand{\Mcal}{\mathcal M}
\newcommand{\dd}{\,\mathrm d}
\newcommand{\e}{\mathrm e}
\title{\textbf{How Random Is the M\"obius Function?\\ Smoothing, Probability, and the Riemann Hypothesis}}
\author{Alberto Verjovsky}
\date{Source: arXiv:2607.25002v2 (CC BY 4.0)}
\begin{document}
\maketitle

\noindent\emph{Source and licence.} \textbf{How Random Is the M\"obius Function?\\ Smoothing, Probability, and the Riemann Hypothesis}. Version arXiv:2607.25002v2 (CC BY 4.0), distributed by arXiv under the Creative Commons Attribution 4.0 International licence, http://creativecommons.org/licenses/by/4.0/, as recorded in the arXiv licence metadata for this exact version and shown on the abstract page https://arxiv.org/abs/2607.25002v2. Full licence notice: \texttt{licenses/arxiv-2607.25002-CC-BY-4.0.txt}.\par\smallskip

\noindent\textit{Editorial scope.} Counted unit: the article's theorem thm:fourier (source0.tex lines 443--458). The article's complete original proof of it is delivered with it (source0.tex lines 1457--1541, one \texttt{proof} environment of 85 lines, which the article places away from the statement). No mathematics is written, repaired or re-derived: the statement and the proof are the article's own lines, in the article's own notation, and no retained line is substituted or re-flowed.  Retained uncounted as the source-local context the proof needs: lines 460--465; lines 1408--1417.  Nothing was omitted from any retained span, so the package carries no omission record.  No retained line contains a citation, so the wrapper carries no bibliography and no bibliography entry.  No retained line contains a figure environment, an image-inclusion command or drawing code, so no image is delivered and the package declares no asset dependency.  Editorial formatting only, in addition: an AMS \texttt{amsart} preamble that restates the article's own declarations and macros used by the retained lines, namely the environments \texttt{theorem}, \texttt{remark} on separate counters, together with the standard packages those lines need.  The retained environments print in this document as Theorem 1, Remark 1 in order of appearance, whereas the article prints the same items with the numbering of its own full document.  The complete native source of the article as distributed in the arXiv e-print for this exact version is bundled unchanged as \texttt{sources/206155/source0.tex}.
\begin{remark}\label{rem:iii-precise}
The set \(Q_\eta\) may depend on \(\eta\).  In the proof of
\textup{(iii)}\(\Rightarrow\)\textup{(i)}, one therefore fixes
\(\eta\) first and then chooses a sufficiently large
\(q\in Q_\eta\).
\end{remark}

\begin{equation}\label{eq:Mertens-local}
 |M(N)|
 \leq C_{q,c}
 \max\left\{
 \sqrt N\,\Mcal_{q,c}(N),
 \,
 N^{\frac12+\frac{1}{2(q+1)}}
 \Mcal_{q,c}(N)^{\frac{q}{q+1}}
 \right\}.
\end{equation}

\begin{theorem}[Local Fourier-moment criterion]\label{thm:fourier}
Fix \(c>0\).  The following assertions are equivalent.
\begin{enumerate}[label=\textup{(\roman*)}]
\item The Riemann hypothesis holds.
\item For every \(\eta>0\) and every finite \(q\ge1\),
\[
\Mcal_{q,c}(N)=O_{\eta,q,c}(N^\eta).
\]
\item For every \(\eta>0\), there is an unbounded set
\(Q_\eta\subseteq[1,\infty)\) such that
\[
\Mcal_{q,c}(N)=O_{\eta,q,c}(N^\eta)
\qquad(q\in Q_\eta).
\]
\end{enumerate}
\end{theorem}

\begin{proof}[Proof of \cref{thm:fourier}]
Assume first that RH holds.  Then for every \(\delta>0\),
\[
 M(x)=O_\delta(x^{1/2+\delta}).
\]
Abel summation (partial summation) gives, for \(t\in\mathbb R\), the
exact identity
\begin{equation}\label{eq:Abel}
 S_N(t)
 =
 M(N)\e^{2\pi iNt}
 -
 2\pi it\int_1^N M(x)\e^{2\pi ixt}\,\dd x.
\end{equation}
Indeed, with \(f(x)=\e^{2\pi ixt}\) and the convention \(M(x)=0\) for
\(x<1\), the general partial summation formula
\[
 \sum_{y<n\leq N}\mu(n)f(n)
 =
 M(N)f(N)-M(y)f(y)-\int_y^N M(x)f'(x)\,\dd x
\]
applied with any \(y\in[0,1)\) gives \(M(y)=0\) and
\(\sum_{y<n\leq N}\mu(n)f(n)=\sum_{n\leq N}\mu(n)f(n)=S_N(t)\), so it
reduces exactly to \eqref{eq:Abel} with no residual boundary term.
Consequently,
\[
 |S_N(t)|
 \ll_\delta
 N^{1/2+\delta}
 +
 |t|\int_1^N x^{1/2+\delta}\,\dd x
 \ll_\delta
 N^{1/2+\delta}(1+N|t|).
\]
For \(|t|\leq c/N\), this gives
\[
 |P_N(t)|\ll_{\delta,c}N^\delta.
\]
Therefore
\[
 \Mcal_{q,c}(N)\ll_{\delta,q,c}N^\delta
\]
for every finite \(q\geq1\).  This proves
\textup{(i)}\(\Rightarrow\)\textup{(ii)}, and
\textup{(ii)}\(\Rightarrow\)\textup{(iii)} is immediate.

Assume now \textup{(iii)}, in the precise form of \cref{rem:iii-precise}:
for every \(\eta>0\) there is an unbounded set \(Q_\eta\subseteq[1,\infty)\)
such that \(\Mcal_{q,c}(N)=O_{\eta,q,c}(N^\eta)\) for every
\(q\in Q_\eta\).

Let \(\varepsilon>0\).  The two choices below must be made in this
order, since the set of admissible exponents \(Q_\eta\) depends on
\(\eta\).  First choose \(\eta>0\) so small that
\[
 \eta<\frac{\varepsilon}{3}.
\]
By hypothesis \(Q_\eta\) is unbounded, so we may then choose
\(q\in Q_\eta\) large enough that
\[
 \frac1{2(q+1)}<\frac{\varepsilon}{3}.
\]
Since \(q\in Q_\eta\), the hypothesis gives
\[
 \Mcal_{q,c}(N)=O_{\eta,q,c}(N^\eta).
\]
The first term in \eqref{eq:Mertens-local} is
\[
 O(N^{1/2+\eta}),
\]
while the second is
\[
 O\left(
 N^{1/2+1/(2(q+1))+\eta q/(q+1)}
 \right)
 =
 O(N^{1/2+\varepsilon}).
\]
Thus
\[
 M(N)=O_\varepsilon(N^{1/2+\varepsilon}).
\]
Since this holds for every \(\varepsilon>0\), the classical Mertens
criterion implies RH.
\end{proof}

\end{document}
